% SSSUP-FONTE-STORICA: fonti-locali/originali/2014-2020.zip
% Anno 2014, esercizio 3. Trasposizione dal sorgente LaTeX storico.
% Verificare editorialmente la trasposizione prima di considerarla definitiva.
\renewcommand{\Anno}{2014}
\renewcommand{\NumeroProblema}{3}

\renewcommand{\Testo}{
Il gioco dell'uomo nero (per due giocatori) può essere descritto nel modo seguente. Si prende un mazzo di carte (comprensivo di un unico jolly), e le si distribuiscono ai due giocatori. Ogni giocatore scarta le coppie di carte di ugual valore (due assi, due re, etc.) che possiede, restando pertanto solo con carte ``spaiate'' (ed uno di loro con il jolly). A questo punto uno dei giocatori pesca a caso una carta dall'altro giocatore: se pesca il jolly lo mette insieme alle altre, se pesca una carta che può appaiare con una delle proprie scarta la coppia che si è formata. Successivamente è l'altro giocatore a pescare una carta, ed il gioco prosegue fino a quando uno dei due giocatori resta con il solo jolly (e l'altro evidentemente senza carte). Il giocatore con il jolly ha perso.

Supponiamo che ad un certo punto del gioco uno dei giocatori (a cui spetta il turno di pescare la carta) sia rimasto con tre carte e l'altro con quattro.

\begin{enumerate}
\item[(a)] Verificare che la probabilità che il gioco non finisca con la vittoria di un giocatore è inferiore a qualsiasi numero reale positivo (ed è pertanto nulla).
\item[(b)] Determinare la probabilità che vinca il giocatore con tre carte.
\item[(c)] Determinare la probabilità di vittoria quando il giocatore cui spetta il turno di pescare possieda $n$ carte ed il suo avversario $n+1$, con $n$ rispettivamente pari e dispari.
\end{enumerate}
}

\renewcommand{\Risposta}{
% SSSUP-BOZZA-AI: Risposta; autore=ChatGPT; data=2026-11-10; stato=DA-VERIFICARE
% SSSUP-PROVENIENZA: ChatGPT (OpenAI), bozza AI 2014 del 2026-11-10; ricostruzione indipendente dal Testo storico; revisione umana pendente.
\begin{enumerate}
\item[(a)] La probabilità che il gioco non termini è $0$.
\item[(b)] La probabilità richiesta è $3/5$.
\item[(c)] La probabilità di vittoria del giocatore di turno
che ha $n$ carte (contro $n+1$) è
\[
 P_n=\begin{cases}
 \dfrac{n+4}{2(n+2)},&n\text{ pari},\\[4pt]
 \dfrac{n+3}{2(n+2)},&n\text{ dispari}.
 \end{cases}
\]
\end{enumerate}
}
\renewcommand{\Aiutino}{
% SSSUP-BOZZA-AI: Aiutino; autore=ChatGPT; data=2026-11-10; stato=DA-VERIFICARE
% SSSUP-PROVENIENZA: ChatGPT (OpenAI), bozza AI 2014 del 2026-11-10; ricostruzione indipendente dal Testo storico; revisione umana pendente.
Dopo lo scarto iniziale, ogni carta ordinaria rimasta in mano
a un giocatore ha la propria compagna in mano all'altro.
Di conseguenza, pescare una carta diversa dal jolly fa
scartare una coppia. Conviene introdurre la probabilità
$P_n$ di vittoria per chi deve pescare e possiede $n$
carte, mentre l'avversario ne possiede $n+1$ (compreso il
jolly), e condizionare alla prima pescata.
}
\renewcommand{\SoluzioneUfficiale}{}
\renewcommand{\SoluzioneRedazionale}{
% SSSUP-BOZZA-AI: SoluzioneRedazionale; autore=ChatGPT; data=2026-11-10; stato=DA-VERIFICARE
% SSSUP-PROVENIENZA: ChatGPT (OpenAI), bozza AI 2014 del 2026-11-10; ricostruzione indipendente dal Testo storico; revisione umana pendente.
Ogni coppia di carte ordinarie dello stesso valore viene
scartata appena si trova nella mano di un giocatore.
Quindi tutte le carte ordinarie ancora in gioco si trovano
in coppie divise fra i due giocatori, una carta per mano.
Se le mani hanno $n$ e $n+1$ carte, il jolly si trova
necessariamente nella mano più numerosa: ciascun giocatore
ha infatti $n$ carte ordinarie.

Indichiamo con $P_n$ la probabilità che vinca chi deve
pescare e ha $n$ carte, mentre l'avversario ne ha $n+1$.
Per $n=0$ il primo giocatore ha già vinto, dunque $P_0=1$.

Alla prima pescata, con probabilità $1/(n+1)$ viene preso
il jolly. Le dimensioni delle mani si scambiano e tocca
all'altro giocatore pescare: la probabilità di vittoria
del giocatore iniziale diventa $1-P_n$.
Con probabilità $n/(n+1)$ viene invece pescata una carta
ordinaria, che forma una coppia e viene scartata. Per
$n\geq2$, al turno successivo l'avversario pesca
necessariamente una carta ordinaria, scartando un'altra
coppia: si ritorna alla medesima situazione con $n-2$
carte nella mano del giocatore iniziale. Il contributo
è pertanto $P_{n-2}$. Se $n=1$, la pescata ordinaria
fa vincere immediatamente il primo giocatore.

Perciò
\[
 P_1=\frac12(1-P_1)+\frac12,
 \qquad P_1=\frac23,
\]
e, per $n\geq2$,
\[
 P_n=\frac1{n+1}(1-P_n)
        +\frac n{n+1}P_{n-2},
 \qquad
 P_n=\frac{1+nP_{n-2}}{n+2}.
\]
In particolare
\[
 P_2=\frac{1+2P_0}{4}=\frac34,
 \qquad
 P_3=\frac{1+3P_1}{5}=\frac35.
\]

Per induzione separata sugli indici pari e dispari,
dalla ricorrenza segue
\[
 P_n=\begin{cases}
 \dfrac{n+4}{2(n+2)},&n\text{ pari},\\[4pt]
 \dfrac{n+3}{2(n+2)},&n\text{ dispari}.
 \end{cases}
\]

Resta da verificare che il gioco finisca quasi certamente.
Finché vi sono $n$ coppie ordinarie, ogni pescata del
giocatore con $n$ carte ha probabilità $n/(n+1)>0$
di eliminare una coppia, mentre la pescata del jolly
si limita a invertire i ruoli dei due giocatori.
La probabilità di pescare il jolly per $k$ tentativi
consecutivi, senza eliminare alcuna coppia, è
$(1/(n+1))^k$, che tende a zero. Quindi, con probabilità
$1$, prima o poi il numero di coppie diminuisce.
Ripetendo l'argomento un numero finito di volte, si
arriva certamente, salvo un evento di probabilità nulla,
alla fine della partita. Questo dimostra (a).
}
\renewcommand{\SoluzioniAlternative}{}
\renewcommand{\Osservazioni}{}
\renewcommand{\Commenti}{}
